\documentclass{article}

\usepackage{amsmath,amssymb}
\usepackage{xurl}
\usepackage[hidelinks]{hyperref}

% Shared commands for notes in docs/paper-gaps/.

\DeclareUnicodeCharacter{2080}{\ifmmode{}_{0}\else\textsubscript{0}\fi}
\DeclareUnicodeCharacter{2081}{\ifmmode{}_{1}\else\textsubscript{1}\fi}
\DeclareUnicodeCharacter{2082}{\ifmmode{}_{2}\else\textsubscript{2}\fi}
\DeclareUnicodeCharacter{2099}{\ifmmode{}_{n}\else\textsubscript{n}\fi}

\newcommand{\Lean}{\textsc{Lean}}
\ExplSyntaxOn
\NewDocumentCommand{\leanid}{m}{
  \ifmmode
    \textsf{\detokenize{#1}}
  \else
    \group_begin:
    \urlstyle{sf}
    \tl_set:Nn \l_tmpa_tl {#1}
    \tl_replace_all:Nnn \l_tmpa_tl {\_} {_}
    \exp_args:NV \path \l_tmpa_tl
    \group_end:
  \fi
}
\ExplSyntaxOff
\providecommand{\tr}{\operatorname{tr}}
\newcommand{\tnleanIssueBase}{https://github.com/LionSR/TNLean/issues/}
\newcommand{\tnleanPrBase}{https://github.com/LionSR/TNLean/pull/}
\newcommand{\tnleanDocBase}{https://sirui-lu.com/TNLean/docs/}
\newcommand{\ghissue}[1]{\href{\tnleanIssueBase#1}{GitHub issue \##1}}
\newcommand{\ghpr}[1]{\href{\tnleanPrBase#1}{PR \##1}}
\newcommand{\doclink}[1]{\href{\tnleanDocBase#1.html}{\path{#1}}}

% Dirac notation and the identity operator, as in blueprint/src/macros/common.tex.
% \providecommand keeps note-local definitions authoritative.
\usepackage{dsfont}
\providecommand{\ket}[1]{|#1\rangle}
\providecommand{\bra}[1]{\langle #1|}
\providecommand{\braket}[2]{\langle #1 | #2 \rangle}
\providecommand{\ketbra}[2]{|#1\rangle\!\langle #2|}
\providecommand{\Id}{\mathds{1}}
\providecommand{\one}{\mathds{1}}
\providecommand{\C}{\mathbb{C}}
\providecommand{\MN}[1]{M_{#1}(\C)}
\providecommand{\MD}{M_D(\C)}

% Verdict marker: \gapnote{<kind>}{<status>}. Kind is the policy
% classification (clarification, local-correction, scope-restriction,
% unfaithful, false-source, open-gap); status is open, wip (elimination
% underway), resolved, or historical. Machine-read by the site index;
% prints a verdict line.
\newcommand{\gapnote}[2]{\par\noindent\textbf{Verdict:} #1 (#2).\par}


% Fallbacks keep this author document compilable when it is built outside its
% source directory and a different command.tex is found first.
\providecommand{\leanid}[1]{\textsf{\detokenize{#1}}}
\providecommand{\doclink}[1]{%
    \href{https://sirui-lu.com/TNLean/docs/#1.html}{\path{#1}}}
\providecommand{\gapnote}[2]{%
    \par\noindent\textbf{Verdict:} #1 (#2).\par}

\title{Domain Walls: Nonzero Tensors, Phases and Blocked Local Actions}
\author{}
\date{2026-09-27}

\begin{document}
\maketitle
\gapnote{local-correction}{resolved}

Ref.~\cite{GarreRubioSchuch2024DomainWall} models the excitations between two
symmetry-broken ground states $A_x$ and $A_y$ of a matrix product unitary
symmetry by domain-wall tensors $e_{xy}$, and describes the symmetry action on
them locally: the action tensors of $g$ on the two regions, applied to $O_g$
acting on the chain $A_x\,e_{xy}\,A_y$, return $B^g_{x,y}$ times the chain
$A_{gx}\,e_{gx,gy}\,A_{gy}$.\footnote{Source traceability:
\path{Papers/2405.00439/MPU-DW.tex}, lines 643--838 (the $\mathbb Z_2$ case,
\texttt{SingleDW}, \texttt{DWopmps}, \texttt{eq:localcdef}) and lines
1889--1933 (the finite-group case, \texttt{eq:localcdefG}, \texttt{PentLB}).}
Two readings of these definitions are adopted as conventions.

\section{Nonzero domain walls and phases}

The source never states that a domain-wall tensor is nonzero, nor that the
scalars $c_{AB}$, $c_{BA}$ and $B^g_{x,y}$ are nonzero. Read literally, the
zero tensor is a domain wall carried to itself with every phase, and a zero
phase is allowed for a nonzero wall. Neither is intended: the tensors model
excitations, and the scalars are called phases and phase factors (lines 712,
1656 and 1894). The formalization requires the walls to be nonzero and the
phases to be nonzero wherever a scalar identity between phases is derived; for
the families of domain walls of Section~IV both conditions are part of the
definition.\footnote{Formal declarations:
\leanid{MPOTensor.GroupFamily.BlockActionData.IsDomainWallAction.mul_eq},
\leanid{MPOTensor.GroupFamily.BlockActionData.IsDomainWallFamily},
\leanid{MPOTensor.GroupFamily.BlockActionData.IsDomainWallAction.mul_eq_lSymbol_of_mul_self_eq_one}.}
Without a nonzero wall no phase is determined, so the literal reading has no
content to formalize.

\section{Blocked local actions}

The source draws the local action at the renormalization fixed point, with
one site of $A_x$ and one site of $A_y$ beside the wall (lines 774--832), and
states that away from the fixed point the same analysis applies after
blocking (line 1358). The formalization states the local action for all words
of $A_x$ and $A_y$ longer than a fixed buffer. The buffered form is needed even
for the fixed-point example of the source: for the CZX symmetry and the
product states $|0\rangle^{\otimes N}$ and $|1\rangle^{\otimes N}$, the wall
$e_{BA}=-|1\rangle$ acquires its phase $c_{BA}=-1$ only with at least one site
to the right of the wall.\footnote{Formal declarations:
\leanid{MPOTensor.GroupFamily.BlockActionData.IsDomainWallAction},
\leanid{CZXCompression.czx_isDomainWallAction_ba}.}

\textbf{Verdict: resolved local fix.} Domain walls are nonzero, their phases
are nonzero, and the local action holds against blocked regions; with these
readings the source's relations $c_{AB}c_{BA}=L_A/L_B=\omega$, \texttt{PentLB}
and \texttt{Intequiv} are proved.

\bibliographystyle{alpha}
\bibliography{references}

\end{document}
